\section{Lean definitions, object-language targets, and completion plan}
\label{sec:formalization}

\subsection{What has actually been compiled}

The archive contains a standalone Lean 4.19.0 project using only the bundled
\texttt{Std} library. Its build was run successfully. The two finite
mathematics modules contain 30 named proved theorems, and their build also
checks three closed examples. The other modules contain compiling
primitive-recursive syntax and representation specifications. There are no
\texttt{sorry}, \texttt{admit}, or new axiom declarations in the source.
Representative axiom reports use \texttt{propext} and, for some results,
\texttt{Quot.sound}; none of those inspected reports uses
\texttt{Classical.choice} or \texttt{sorryAx}. The reports are included.

These are semantic Lean proofs of finite algebra. They are not described as
completed object-iPRA derivations of the analytic results. ProveIt's existing
object calculus uses a separately pinned toolchain; its imported sources were
inspected but not rebuilt for this package. The standalone project deliberately
has no hidden dependency on that checkout.

\begin{longtable}{>{\raggedright\arraybackslash}p{.26\textwidth}>{\raggedright\arraybackslash}p{.65\textwidth}}
\toprule
Compiled file & Verified content\\
\midrule\endhead
\texttt{Finite.lean} & Bounded sums, finite event mass, complement and disjoint
additivity, finite union bounds, product rectangles, Markov, a centered
power-Markov form of Chebyshev, and stability of event mass in symmetric
difference.\\
\texttt{L1Finite.lean} & Signed finite sums, the absolute-sum inequality,
expectation stability, and the finite $L^1$ triangle inequality.\\
\texttt{PRSyntax.lean} & Arity-indexed PR syntax, finite evaluation derivations,
and a connection between one fixed syntax tree and its semantic function.
No universal PR evaluator is defined.\\
\texttt{Specifications.lean} & Positive-denominator rational codes, fixed-PR
rational sequences, regular real specifications, PR event tables, common
prefix refinement, and a regular event specification.\\
\bottomrule
\end{longtable}

The compiled regular-real and event specifications use the stronger
one-sided convention $d(x_k,x_l)\le2^{-k}$ for $l\ge k$.
The report's symmetric convention is $d(x_k,x_l)\le2^{-k}+2^{-l}$.
A one-bit shift converts a symmetric name to the one-sided convention, and
the latter already implies the symmetric one. This harmless reindexing is
explicit and should be retained when connecting the interfaces.

\subsection{All representation data for the continuous layer}

The following additional structures are specifications to implement; they
are not postulated instances. Their arithmetic contents are completely
determined by the definitions in the report.

\begin{longtable}{>{\raggedright\arraybackslash}p{.22\textwidth}>{\raggedright\arraybackslash}p{.69\textwidth}}
\toprule
Interface & Fields and exact obligation\\
\midrule\endhead
\texttt{PRTable} & A fixed PR depth program $\ell(k)$ and three fixed PR
programs on $(k,i)$ for positive/negative numerator and denominator minus
one. Only $i<2^{\ell(k)}$ is read. Prefix lifting uses
$i\mapsto\lfloor i/2^{D-\ell(k)}\rfloor$.\\
\texttt{RandomName0} & A \texttt{PRTable} and a derivation of the all-precision
$d_0$ inequality, where $d_0$ is the finite rational fold of
$\min(1,|x-y|)$.\\
\texttt{RandomName1} & The same, using the fold of $|x-y|$. Expectation is
the rational sum divided by the row count.\\
\texttt{RandomName2} & The same, with
$I((x_k-x_l)^2)\le(2^{-k}+2^{-l})^2$. No square-root test is needed
in the validity predicate.\\
\texttt{EventName} & Boolean tables and symmetric-difference counts;
the cleared formula is given in Section~\ref{sec:scope}.\\
\texttt{BandSpec} & Fixed PR $b(k),R(k)$ and proofs that for every
$r\ge R(k)$ the finite band count at width $2^{-b(k)}$ is at most $2^{-k}$.\\
\texttt{CompactFunction} & A fixed rational evaluator, local bound, continuity
modulus, and the evaluator's finite consistency and oscillation proofs.
Optional increment and derivative contracts support calculus.\\
\texttt{KernelSpec} & A nonnegative \texttt{CompactFunction} on each compact
range, its listed breakpoints, mass-tail radius, and a fixed positive
normalization witness (K1)--(K2).\\
\texttt{MomentSpec} & An integer order $p$, PR cutoff, and finite annular-tail
proofs for $|x|^p w(x)$. A numerical bound on the moment is additional data
when a theorem, such as the CLT, requests one.\\
\texttt{ProductFamily} & A fixed PR bit-position map, injectivity/disjoint-row
counting proofs, and uniform finite-dimensional approximation schedules.
These implement the independent copies used here.\\
\texttt{SmoothTest} & PR evaluators for $h,h',h''$, bounds, and a proved
Taylor remainder $|R(x,u)|\le L|u|^3/6$, with the finite evaluation errors
included. The quintic cutoff supplies a concrete instance.\\
\texttt{TailUnionSpec} & A jointly PR event family and a PR rate for distances
between its finite unions. A summable mass-tail provider is one constructor.\\
\bottomrule
\end{longtable}

For clarity, the compiled syntax begins with the following shape (the source
contains its fully typed Lean form):
\begin{lstlisting}
inductive PR : Nat -> Type where
  | zero (n : Nat) : PR n
  | succ : PR 1
  | proj (n : Nat) (i : Fin n) : PR n
  | comp (outer : PR m) (inner : Fin m -> PR n) : PR n
  | primrec (base : PR n) (step : PR (n + 2)) : PR (n + 1)
\end{lstlisting}
A finite \texttt{Eval} derivation has separate constructors for the defining
equations. \texttt{FixedPRFunction} carries one code, its semantic function,
and evaluation-graph proofs for that fixed code. The function field is a
metatheoretic convenience, not an admission of all Lean functions into iPRA.
The final object version must carry quoted PR symbols and actual proofs of
the numerical contracts in \texttt{Prf PRAAx}.

\subsection{The main embedded theorem targets}

For a fixed table program $X$ let $\operatorname{Reg}_0(X)$ denote the
all-precision first-order formula above, and similarly for the other
metrics. Let $\operatorname{Band}(X,t)$ abbreviate the quantified finite
band test with its fixed $b,R$ programs. The intended targets are genuine
formulas over natural codes, such as
\begin{align*}
 &\operatorname{Reg}_1(X)\ \longrightarrow\
   \operatorname{Reg}_{\mathbb R}(\operatorname{Expect}(X)),\\
 &\operatorname{Reg}_0(X)\wedge\operatorname{Band}(X,t)
   \ \longrightarrow\
   \operatorname{Reg}_{\rm ev}(\operatorname{Threshold}(X,t)),\\
 &\operatorname{Kernel}(w)\ \longrightarrow\
   \operatorname{Reg}_0(\operatorname{HistogramSampler}(w)),\\
 &\operatorname{Kernel}(w)\wedge\operatorname{Moment}_2(w)
   \ \longrightarrow\
   \operatorname{Reg}_2(\operatorname{ReindexSampler}(w)).
\end{align*}
Each named construction on the right is a fixed PR program built from the
specified input definitions. There is no free functional quantifier in these
schemas. Event and random-variable equality are all-precision relations;
they are not tested by an evaluator.

For the strong law, quote the finite sum, maximum, event indicator, product
weight, and error-allowance programs. The endpoint is the cleared version of
\eqref{eq:strong-window}, universally quantified over $N,M$ and rational
precision parameters under the fixed presentation contracts. Its derived
union name uses $T\ge4V\,2^k/\varepsilon^2$ for residual precision
$2^{-k}$. For the CLT, quote the quintic cutoff and the finite Taylor-error
test; the target is a pair of expectation inequalities whose error is
\eqref{eq:cdf-clt}, or its squared version with all positive square roots
eliminated by rational intervals. The sample-size condition
\eqref{eq:clt-rate} is already integral arithmetic.

In ProveIt the final endpoint has the form
\begin{lstlisting}
-- Desired proof type after the stated formula has been constructed:
Prf PRAAx ([] : List (Fm PTm 0)) phi
\end{lstlisting}
This is a target type, not an asserted theorem constant in the delivered
code. A law field in a structure, an axiom, a standard-model truth proof,
or a program that merely checks examples would not fill this endpoint.

\subsection{A finite dependency plan with explicit induction invariants}

\begin{longtable}{>{\raggedright\arraybackslash}p{.08\textwidth}>{\raggedright\arraybackslash}p{.38\textwidth}>{\raggedright\arraybackslash}p{.45\textwidth}}
\toprule
Stage & Construction and proof & Quantifier-free invariant or finite bridge\\
\midrule\endhead
1 & Quote rational codes, list access, refinement, sums and norms. &
Valid denominators, running sum identity/inequality, and prefix projection
composition. Reuse the existing signed and bounded-sum arithmetic.\\
2 & Prove finite probability and metric estimates. &
Weighted partial-sum bounds; rowwise indicator identities; expanded
nonnegative squares. Existing \texttt{Expectation} providers already
cover much of this layer.\\
3 & Reify regularity and extensional equality. &
Instantiate input contracts at requested finite indices, use triangle
inequalities, then remove rational slack. Do not induct on a quantified
regularity predicate.\\
4 & Quote compact quadrature and calculus. &
Accumulated oscillation/error budget; telescoping increment identity;
partial Taylor sum with its geometric tail.\\
5 & Normalize kernels and realize laws. &
Nonnegative dyadic masses, sum at most one, residual at zero, monotone
cumulative allocation, finite clipped coupling and moment-tail bounds.\\
6 & Extend events, products and conditioning. &
Symmetric-difference budget, bounded-multiplier truncation bound, finite
double-sum exchange, positive denominator guard.\\
7 & Lift weak and maximal inequalities. &
First-crossing partition and centered product identities; squared error
and boundary-band budgets before the limit. Reuse \texttt{WeakLaw},
\texttt{FirstCrossing}, and the explicitly scoped \texttt{Kolmogorov*} providers.\\
8 & Construct strong-law exceptional events. &
Finite dyadic-block geometric sum, finite-union residual estimate,
then explicit diagonal schedule for the countable precision scales.\\
9 & Prove the quantitative CLT. &
After $j$ replacements, the expectation discrepancy is bounded by the
sum of the first $j$ Taylor/moment errors. Verify the polynomial cutoff
identities and allocate the finite moment-rounding errors.\\
10 & Export and audit. &
Construct the actual closed \texttt{Prf PRAAx} endpoints, then use the
existing transport theorem to the chosen concise presentation. Check
axiom and dependency reports at those endpoints.\\
\bottomrule
\end{longtable}

Stages 1--2 have substantial existing object-level support, plus the new
standalone finite checks. The new continuous constructions in stages 3--9
are supplied here as human proofs and formalization specifications. Their
full embedded implementation remains work to do. The finite examples and
compiled specification records do not turn that remaining work into a
completed formal proof.

\clearpage
\section{Coverage, evidence, and explicit limits}
\label{sec:coverage}

\begin{longtable}{>{\raggedright\arraybackslash}p{.32\textwidth}>{\raggedright\arraybackslash}p{.58\textwidth}}
\toprule
Topic & Result and exact scope\\
\midrule\endhead
Probability space and events & Coherent rational cylinders, represented
measure algebra, finite probability axioms, countable additivity for
supplied finite-union tail moduli.\\
Random variables and expectation & $L^0$, $L^1$, $L^2$ names; law through
bounded observables; expectation, covariance, product and multiplier
closure under their stated integrability assumptions.\\
Continuous distributions & Uniform, exponential, normal, Cauchy, and Cantor
singular continuous law; explicit tails, moments when finite, CDFs,
histogram realization, and separated pointwise quantiles.\\
Joint laws and conditioning & Explicit independent copies, represented
Fubini with absolute tails, convolution, finite conditional averaging,
Bayes updates with positive evidence, Gaussian conditional density.\\
Convergence & PR-modulated completion, quantitative Vitali and dominated
convergence, canonical dyadic-martingale convergence, certified first
Borel--Cantelli fragment.\\
Laws of large numbers & Finite weak-law bound, finite-window maximal/strong
bound, and effective exceptional-event certificates when threshold events
have the required presentation.\\
Central limit theorem & Smooth Lindeberg replacement and uniform CDF error
with explicit $O(\delta^{-8})$ sufficient sample size under a third-moment
bound. An optimal Berry--Esseen estimate is not claimed.\\
Stochastic processes & Finite-horizon continuous-state transitions, bounded
stopping, and an explicit autoregressive Gaussian stationary-law example.
No general continuous-time path or stochastic-integral construction.\\
\bottomrule
\end{longtable}

The Python reference uses integer and rational arithmetic in all acceptance
checks. The recorded run passes 2,963 checks: cylinder refinements and
moments, conditional averaging and tower identities, independence,
clipping, finite first-crossing inequalities, cutoff-polynomial identities,
and certified Gaussian enclosures. In particular it certifies
\[
 \frac{882141}{1048576}\le\Phi(1)\le\frac{55143}{65536},
 \qquad 0.841<\Phi(1)<0.842.
\]
The decimal endpoints on the right abbreviate exact rationals. The narrower
fraction interval is obtained by positive Taylor sums, inversion, midpoint
quadrature, and the Gaussian tail inequality. Floating-point libraries are
not proof oracles in this calculation. Such finite checks validate the
implementation and constants; the general conclusions rest on the human
proofs and, for the finite Lean layer, kernel checking.

For perspective, $V=1$, $\varepsilon=1/10$, $\delta=1/20$ give sufficient
sample sizes $2{,}000$ for the weak law and $8{,}000$ for the uniform
strong-law window. The elementary CLT with $B=2$ and desired error $1/10$
gives the much larger sufficient value $40{,}960{,}000{,}000{,}000$.
That conservative constant is a cost of the elementary smoothing proof,
not a statement about the actual convergence speed of a particular law.
The package's purpose is a transparent proof-theoretic foundation; sharper
certified algorithms can replace these providers without changing the
representation interfaces.

The precise exclusions are as informative as the positive results.
Arbitrary measurable-set comprehension, regular probabilities for all
effective Borel descriptions, general rate extraction from qualitative
convergence, and unrestricted disintegration are not available for free.
Their absence is supported by the explicit halting obstructions and the
conditioning results cited above, rather than by a claim that continuous
mathematics itself is impossible in iPRA. The continuous mathematics succeeds
when the finite information consumed by each operation is stated and proved.
